\seckiny
\section{Introduction}
\seckiny

Deep reinforcement learning has recently enjoyed successes in many domains~\cite{mnih-dqn-2015,schulman2015trust,levine2015endtoend,mnih2016asynchronous,lillicrap2015continuous}. 
%
Nevertheless, long-term credit assignment remains a major challenge for these methods, especially in environments with sparse reward signals, such as the infamous Montezuma's Revenge ATARI game. 
It is symptomatic that the standard approach on the ATARI benchmark suite~\cite{bellemare-ale} is to use an action-repeat heuristic, where each action translates into several (usually $4$) consecutive actions in the environment. 
Yet another dimension of complexity is seen in non-Markovian environments that require memory -- these are particularly challenging, since the agent has to learn which parts of experience to store for later, using only a sparse reward signal. 


%%
The framework we propose takes inspiration from feudal reinforcement learning (FRL) introduced by \citet{dayan1993feudal}, where levels of hierarchy within an agent communicate via explicit goals. 
Some key insights from FRL are that goals can be generated in a top-down fashion, and that goal setting can be decoupled from  goal achievement; a level in the hierarchy communicates to the level below it what must be achieved, but does not specify \emph{how} to do so.  Making higher levels reason at a lower temporal resolution naturally structures the agents behaviour into temporally extended sub-policies.  

The architecture explored in this work is a fully-differentiable neural network with two levels of hierarchy (though there are obvious generalisations to deeper hierarchies).
The top level, the “Manager”, sets goals at a lower temporal resolution in a latent state-space that is itself \textit{learnt} by the Manager.
The lower level, the “Worker”, operates at a higher temporal resolution and produces primitive actions, conditioned on the goals it receives from the Manager. The Worker is motivated to follow the goals by an intrinsic reward.
However, significantly, no gradients are propagated between Worker and Manager; the Manager receives its learning signal from the environment alone. In other words, the Manager learns to select latent goals that maximise extrinsic reward.

%
The key contributions of our proposal are: (1) A consistent, end-to-end differentiable model that embodies and generalizes the principles of FRL. (2) A novel, approximate \textit{transition policy gradient} update for training the Manager, which exploits the semantic meaning of the goals it produces. (3) The use of goals that are directional rather than absolute in nature. (4) A novel RNN design for the Manager -- a dilated LSTM -- which extends the longevity of the recurrent state memories and allows gradients to flow through large hops in time, enabling effective back-propagation through hundreds of steps.
%

Our ablative analysis (Section \ref{subsec:ablative}) confirms that transitional policy gradient and directional goals are crucial for best performance. Our experiments on a selection of ATARI games (including the infamous Montezuma's revenge) and on several memory tasks in the 3D DeepMind Lab environment~\cite{beattie2016labyrinth} show that FuN significantly improves long-term credit assignment and memorisation. 
